\documentclass[10pt,a4paper]{amsart}
\usepackage[margin=19mm]{geometry}
\usepackage{amsmath,amssymb,mathtools}
\usepackage[scr=rsfs,cal=euler]{mathalfa}
\usepackage{xfrac}
\usepackage[unicode,hidelinks,bookmarks=false]{hyperref}
\newcommand{\ie}{i.e.\ }
\newcommand{\cf}{cf.\ }
\newcommand{\resp}{respectively\ }
\newcommand{\Subsection}{\subsection}
\providecommand{\nobreakdash}{\nobreak\hbox{-}}
\setlength{\emergencystretch}{2em}
\multlinegap0pt
\usepackage{bm}
\usepackage{bbm}
\usepackage{dsfont}
\usepackage{xspace}
\usepackage{cancel}
\usepackage{mathtools}
\usepackage{amsthm}
\makeatletter
\@ifundefined{thm}{\newtheorem{thm}{Thm}}{}
\@ifundefined{lem}{\newtheorem{lem}[thm]{Lemma}}{}
\makeatother
\DeclareMathOperator{\Tr}{Tr}
\newcommand{\C}{\mathbb{C}}
\newcommand{\id}{\mathrm{id}}
\title[A quantum Frucht's theorem and quantum automorphisms of quan]{A quantum Frucht's theorem and quantum automorphisms of quantum Cayley graphs (Lemma 6.5)}
\author{Michael Brannan, Daniel Gromada, Junichiro Matsuda, Adam Skalski, Mateusz Wasilewski}
\date{Source: arXiv:2503.11149v2 (CC BY 4.0)}
\begin{document}
\maketitle

\noindent\emph{Source and licence.} A quantum Frucht's theorem and quantum automorphisms of quantum Cayley graphs. Version arXiv:2503.11149v2 (CC BY 4.0), distributed under the Creative Commons Attribution 4.0 International licence, as recorded in the arXiv licence metadata for this exact version and shown on the abstract page https://arxiv.org/abs/2503.11149v2. Full rights notice: licenses/arxiv-2503.11149-CC-BY-4.0.txt.\par\smallskip

\noindent\textit{Editorial scope.} Counted unit: the Lem printed as Lemma 6.5 of the article (source0.tex lines 1469--1475, source label \texttt{PeterWeyl}), delivered together with the article's own complete original proof of it (source0.tex lines 1476--1530, one \texttt{proof} environment of 55 lines). No mathematics is written, repaired or re-derived: statement and proof are the article's own text line for line. Required context retained uncounted: none. Every definition, display and label of those lines is retained unchanged. Omitted from this package and not redistributed: 1--1468, 1531--1798. Nothing inside a retained excerpt is omitted or altered: every retained body is delivered exactly as the article's lines are written, so no omission receipt of any kind accompanies this package. Citations: the retained excerpts carry no citation command at all, so no bibliography entry of the article is transcribed and no reference is redistributed. Reference adaptation: none was needed and none was made; every reference inside the delivered unit resolves inside it, so no unresolved reference or citation is produced. Printed slips kept literal: no printed word, symbol, hypothesis or number of the counted statement or of its proof was altered. Layout and class: the article's own class is \texttt{amsart} with packages including inputenc, biblatex, amsmath, amstext, amsfonts, amssymb, dsfont, amsthm, bbm, bm, mathtools, csquotes, color, comment; the wrapper is the lane's pinned \texttt{amsart} editorial wrapper at 10pt on A4, which restates the theorem-like environments the retained text needs and adds the article's own macro definitions that the retained text needs (\texttt{\textbackslash C}, \texttt{\textbackslash Tr}, \texttt{\textbackslash id}), with extra wrapper packages (bm, bbm, dsfont, xspace, cancel, mathtools). The delivered numbering is the unit's own. Assets: no retained span contains a figure environment, a graphics-inclusion command, a PDF-page inclusion, a table or a TikZ picture; no image file is copied into this package, the package declares no asset dependency and no image file is redistributed with it. Licence: version arXiv:2503.11149v2 of this article is distributed under the Creative Commons Attribution 4.0 International licence, as recorded in the arXiv licence metadata for this exact version and shown on the abstract page https://arxiv.org/abs/2503.11149v2. A full rights notice is delivered at licenses/arxiv-2503.11149-CC-BY-4.0.txt. Characters: every retained excerpt is pure ASCII, so the delivered engine, XeLaTeX with its default Latin Modern text font, reports no missing character.
\begin{lem} \label{PeterWeyl}
Consider the  isomorphism $\rho:\C[\Gamma]\to \bigoplus_{\pi \in \textup{Irr}(\Gamma)} M_{n_{\pi}}$, given by the (inverse Fourier) identification $\lambda_g \mapsto \bigoplus_{\pi \in \textup{Irr}(\Gamma)} \pi(g)$. Fix $\pi \in \textup{Irr}(\Gamma)$ and vectors $\xi,\eta \in H_\pi$, and set $x=|\xi\rangle \langle \eta | \in M_{n_{\pi}}$. Then 
\[ \rho^{-1} (x) = \frac{ n_\pi}{|\Gamma|}\sum_{g\in \Gamma}
\langle \pi(g) \eta,  \xi \rangle\,\lambda_g.\]
Further if $P=\id_{H_\pi}$ then we have
\[ \rho^{-1} (P) = \frac{ n_\pi}{|\Gamma|}\sum_{g\in \Gamma} \Tr(\pi(g)^*) \, \lambda_g.\] 
\end{lem}

\begin{proof}
Set $n=n_\pi$ and fix an orthonormal basis $(e_1, \ldots, e_n)$ of $H_\pi$.
By linearity it suffices to prove the result for $\xi=e_i, \eta = e_j$ for some $i,j=1,\ldots,n$. Set $y=\frac{ n_\pi}{|\Gamma|}\sum_{g\in \Gamma}
\langle \pi(g) e_j,  e_i \rangle\,\lambda_g$. We need to check that $\rho(y) =x$, so that
\[ \rho(y)_\pi = |e_i\rangle \langle e_j|,\]
\[ \rho(y)_{\pi'} =0, \;\;\; \textup{for} \;\pi' \in \textup{Irr}(\Gamma), \pi'\neq \pi.\]
But the two formulas displayed above are equivalent to the following:
\[ \frac{ n_\pi}{|\Gamma|}\sum_{g \in \Gamma}
\overline{ \langle e_i,  \pi(g) e_j \rangle} \langle e_k,  \pi(g) e_l \rangle = \langle e_k, e_i \rangle \langle e_j, e_l \rangle,    \]
\[ \frac{ n_\pi}{|\Gamma|}\sum_{g\in \Gamma}
\overline{ \langle e_i,  \pi(g) e_j \rangle} \langle \zeta,  \pi'(g) \kappa \rangle = 0  \]
for all $k,l=1,\ldots,n_\pi$ and all vectors $\zeta, \kappa \in H_{\pi'}$, where $\pi' \in \textup{Irr}(\Gamma), \pi'\neq \pi$. These are however precisely the Peter-Weyl orthogonality relations.

The second formula follows from the first one by linearity.


%{\color{red}
% \[
%     \rho^{-1}\left(\bigoplus_{\pi\in\Irr(\Gamma)} x_\pi\right)
%     = \frac{1}{|\Gamma|} \sum_{g\in \Gamma} \sum_{\pi\in\Irr(\Gamma)} n_\pi \Tr(\pi(g)^* x_\pi)\lambda_{g}.
% \]
% Note that $\lambda\cong\bigoplus_\pi \pi^{\oplus n_\pi}$
% by Peter-Weyl theorem. For $g\in \Gamma$, we have
% \begin{align*}
%     \rho^{-1}\rho(g) 
%     &= \frac{1}{|\Gamma|} \sum_{g'\in \Gamma} \sum_{\pi\in\Irr(\Gamma)} n_\pi \Tr(\pi(g')^* \pi(g)) \lambda_{g'}
%     = \frac{1}{|\Gamma|} \sum_{g'\in \Gamma} \sum_{\pi\in\Irr(\Gamma)} \Tr(\pi^{\oplus n_\pi}({g'}^{-1}g)) \lambda_{g'}
%     \\
%     &= \frac{1}{|\Gamma|} \sum_{g'\in \Gamma}  \Tr(\bigoplus_{\pi\in\Irr(\Gamma)}\pi^{\oplus n_\pi}({g'}^{-1}g)) \lambda_{g'}
%     = \sum_{g'\in \Gamma} \frac{1}{|\Gamma|} \Tr(\lambda_{{g'}^{-1}g})\lambda_{g'}
%     = \sum_{g'\in \Gamma} \delta_{g',g}\lambda_{g'}=\lambda_{g}.
% \end{align*}
% Thus for $x=|\xi\rangle \langle \eta|$ and $P=\id_{H_\pi}$,
%\begin{align*}
%     \rho^{-1}(x)&= \frac{1}{|\Gamma|} \sum_{g\in \Gamma} n_\pi \Tr(\pi(g)^* |\xi\rangle \langle \eta|)\lambda_{g}
%     =\frac{n_\pi}{|\Gamma|}\sum_{g\in \Gamma}  \langle{\pi(g) \eta, \xi}\rangle \lambda_{g},
%     \\
    % \rho^{-1}(P)&=\frac{n_\pi}{|\Gamma|}\sum_{g\in \Gamma}  \Tr(\pi(g)^*) \lambda_{g}
    % =\frac{n_\pi}{|\Gamma|}\sum_{g\in \Gamma}  \ol{\Tr(\pi(g))} \lambda_{g}.
    % \\
%    \sum_{i=i}^{n_\pi}\langle{\pi(g) e_i, e_i}\rangle
 %   &= \sum_{i=i}^{n_\pi}\Tr(|e_i\rangle \langle \pi(g) e_i|)
%    = \Tr(\sum_{i=i}^{n_\pi}|e_i\rangle \langle  e_i|\pi(g)^*)
%    = \Tr(\id_{H_\pi}\pi(g)^*) = \Tr(\pi(g)^*).
%\end{align*}
% Just in case:
% \begin{align*}
%     \rho\rho^{-1}(\bigoplus_\pi x_\pi)
%     &=\frac{1}{|\Gamma|}\sum_{g\in \Gamma}  \sum_{\pi\in\Irr(\Gamma)} n_\pi \Tr(\pi(g)^* x_\pi)\rho(\lambda_g)
%     = \frac{1}{|\Gamma|}\sum_{g\in \Gamma} \Tr(\lambda_g^* \bigoplus_\pi x_\pi^{\oplus n_\pi}) \rho(\lambda_g)
%     \\
%     &=\rho(\bigoplus_\pi x_\pi^{\oplus n_\pi})=\bigoplus_\pi x_\pi
% \end{align*}
%}
\end{proof}

\end{document}
